\section{Introducción}
\label{sec:introduccion}

Contexto del problema y motivación. Las citas se escriben con \verb|\textcite| cuando el
autor forma parte de la oración, como propone \textcite{codd1970}, o con \verb|\parencite|
al final de la idea \parencite{han2011, elmasri2016}. Los acrónimos se definen una sola vez
en \texttt{acronimos.tex}: la primera aparición de \ac{KDD} se expande y las siguientes
(\ac{KDD}) muestran solo la sigla.

\subsection{Objetivos}

\begin{itemize}
  \item \textbf{General:} qué se busca lograr con el trabajo.
  \item \textbf{Específicos:}
  \begin{enumerate}
    \item Diseñar el modelo de datos (\ac{DER}) del dominio.
    \item Implementar el esquema en un \ac{SGBDR} (MySQL 8.4).
    \item Analizar los datos con consultas \ac{SQL} y técnicas de minería de datos.
  \end{enumerate}
\end{itemize}

\subsection{Organización del documento}

La \cref{sec:marco} resume los conceptos teóricos, la \cref{sec:metodologia} describe los
datos y el procedimiento, la \cref{sec:resultados} presenta los resultados y la
\cref{sec:conclusiones} cierra con las conclusiones.

\pendiente{Ejemplo de nota que solo se ve con la opción \texttt{borrador}.}
